\begin{abstract}
Compilers that target machines with physically distributed memories must decide where data lives, where computation runs, and what communication connects the two. These decisions were posed as compiler problems in the late 1980s for hypercubes and multicomputers, standardised in High Performance Fortran in the 1990s, largely left to hand-written message passing in the 2000s, and have returned in the 2020s as sharding propagation for accelerator meshes and wafer-scale machines. This article surveys that history with a reproducible method: an operational definition of a distributed-memory compiler, an openly maintained and schema-validated catalog, and a taxonomy on independent axes that separates the decision procedure for decomposition (directives, closed-form analysis, analytic optimisation, or profiling and search) from the representation of ownership and from the communication family. Coding the 54 catalog papers whose full text was read, we test seven generalizations across the eras. A run-time stage accompanies data-dependent ownership in every read system that describes it. The cost models moved from analytic network models in the 1990s to measured or simulated programs from 2019. The leading optimisation objective moved from message count to memory fit and productivity. Block-cyclic and mapping-array decompositions disappear after about 2000. Pure element-level communication disappears after the 2000s, and owner-computes survives in kernel compilers while sharding systems replace it with propagation. We state the evidence strength and the refuting observation for each generalization, and close with the open problems they leave, including the comparison of measured and modelled cost that no read paper makes.
\end{abstract}
